% AUTOMATISCH ERZEUGT aus dem SysML-Modell mit tools/generate_docs.py.
% Nicht von Hand aendern: Inhalte im Modell pflegen, Layout im Kapitel.
\begin{tikzpicture}
% Layout des Zustandsdiagramms (Betriebszustaende). Von Hand in Overleaf pflegen.
% Inhalt (Zustaende, Uebergaenge, Ereignisse) kommt aus dem Modell:
% tables/generated/state_machine_a1.tex, erzeugt mit tools/generate_docs.py.
% Hier stehen nur:
%  - Stile,
%  - je eine Koordinate pos-<Zustand> pro Zustand (fehlt eine, meldet der
%    Generator einen Fehler),
%  - optional ein Stil e-<von>-<nach> je Pfeil, z. B. fuer Boegen oder die
%    Lage der Beschriftung.
\tikzset{
  state/.style={draw, rounded corners=3pt, align=center, minimum width=26mm,
    minimum height=7mm, fill=blue!5, font=\footnotesize\ttfamily},
  terminal/.style={state, fill=red!8},
  transition/.style={-{Latex[length=1.8mm]}, semithick},
  event/.style={font=\scriptsize, fill=white, fill opacity=0.85, text opacity=1,
    inner sep=1pt},
  note/.style={draw, dashed, align=left, font=\scriptsize, text width=62mm,
    fill=yellow!10, inner sep=4pt},
  % Einzelne Pfeile
  e-completed-ready/.style={out=110, in=70, looseness=0.6},
  e-blocked-ready/.style={out=150, in=55, looseness=0.8},
  e-ready-rejected/.style={bend left=40},
  e-rejected-ready/.style={bend left=40},
  e-driving-avoidingObstacle/.style={bend left=40},
  e-avoidingObstacle-driving/.style={bend left=40},
  e-paused-orderAccepted/.style={bend left=35},
  e-orderAccepted-paused/.style={bend left=35},
  e-chargingInterfaceReached-ready/.style={bend left=45},
  e-chargingInterfaceReached-orderAccepted/.style={out=40, in=-110,
    every node/.append style={pos=0.3}},
  e-fault-ready/.style={out=30, in=-150},
  e-safeState-ready/.style={out=60, in=-140},
}
\coordinate (pos-off) at (0, 0);
\coordinate (pos-initializing) at (5, 0);
\coordinate (pos-ready) at (10.5, 0);
\coordinate (pos-orderAccepted) at (17, 0);
\coordinate (pos-driving) at (23, 0);
\coordinate (pos-completed) at (29, 0);
\coordinate (pos-rejected) at (10.5, 3.6);
\coordinate (pos-avoidingObstacle) at (23, 3.6);
\coordinate (pos-blocked) at (29, 3.6);
\coordinate (pos-stopped) at (10.5, -3.8);
\coordinate (pos-paused) at (17, -3.8);
\coordinate (pos-fault) at (3, -3.8);
\coordinate (pos-safeState) at (3, -7.6);
\coordinate (pos-chargingInterfaceReached) at (10.5, -7.6);
\coordinate (pos-returningToChargingStation) at (18.5, -7.6);
\coordinate (pos-returnPaused) at (25, -11);
\coordinate (pos-note) at (26.5, -3.8);

\node[state] (s-off) at (pos-off) {off};
\node[state] (s-initializing) at (pos-initializing) {initializing};
\node[state] (s-ready) at (pos-ready) {ready};
\node[state] (s-orderAccepted) at (pos-orderAccepted) {orderAccepted};
\node[state] (s-rejected) at (pos-rejected) {rejected};
\node[state] (s-driving) at (pos-driving) {driving};
\node[state] (s-avoidingObstacle) at (pos-avoidingObstacle) {avoidingObstacle};
\node[state] (s-paused) at (pos-paused) {paused};
\node[state] (s-blocked) at (pos-blocked) {blocked};
\node[state] (s-completed) at (pos-completed) {completed};
\node[state] (s-stopped) at (pos-stopped) {stopped};
\node[state] (s-returningToChargingStation) at (pos-returningToChargingStation) {returningToChargingStation};
\node[state] (s-returnPaused) at (pos-returnPaused) {returnPaused};
\node[state] (s-chargingInterfaceReached) at (pos-chargingInterfaceReached) {chargingInterfaceReached};
\node[terminal] (s-fault) at (pos-fault) {fault};
\node[terminal] (s-safeState) at (pos-safeState) {safeState};
\fill (s-off.west) ++(-6mm,0) circle (1.2mm) coordinate (start);
\draw[transition] (start) -- (s-off);
\draw[transition, e-off-initializing/.try] (s-off) to node[event, align=center] {PowerOn} (s-initializing);
\draw[transition, e-initializing-ready/.try] (s-initializing) to node[event, align=center] {InitializationComplete} (s-ready);
\draw[transition, e-ready-orderAccepted/.try] (s-ready) to node[event, align=center] {OrderValidated} (s-orderAccepted);
\draw[transition, bend left=14, e-ready-rejected/.try] (s-ready) to node[event, align=center] {OrderRejected} (s-rejected);
\draw[transition, e-orderAccepted-driving/.try] (s-orderAccepted) to node[event, align=center] {RouteReady} (s-driving);
\draw[transition, e-driving-completed/.try] (s-driving) to node[event, align=center] {TargetReached} (s-completed);
\draw[transition, e-driving-blocked/.try] (s-driving) to node[event, align=center] {RouteBlocked} (s-blocked);
\draw[transition, e-driving-paused/.try] (s-driving) to node[event, align=center] {OperationalStop} (s-paused);
\draw[transition, bend left=14, e-orderAccepted-paused/.try] (s-orderAccepted) to node[event, align=center] {OperationalStop} (s-paused);
\draw[transition, e-avoidingObstacle-paused/.try] (s-avoidingObstacle) to node[event, align=center] {OperationalStop} (s-paused);
\draw[transition, bend left=14, e-paused-orderAccepted/.try] (s-paused) to node[event, align=center] {Resume [ResumePermitted]} (s-orderAccepted);
\draw[transition, e-ready-returningToChargingStation/.try] (s-ready) to node[event, align=center] {Return} (s-returningToChargingStation);
\draw[transition, e-paused-returningToChargingStation/.try] (s-paused) to node[event, align=center] {Return} (s-returningToChargingStation);
\draw[transition, e-chargingInterfaceReached-orderAccepted/.try] (s-chargingInterfaceReached) to node[event, align=center] {Resume [ResumePermitted]} (s-orderAccepted);
\draw[transition, e-returningToChargingStation-chargingInterfaceReached/.try] (s-returningToChargingStation) to node[event, align=center] {ChargingAreaReached} (s-chargingInterfaceReached);
\draw[transition, bend left=14, e-returningToChargingStation-returnPaused/.try] (s-returningToChargingStation) to node[event, align=center] {OperationalStop} (s-returnPaused);
\draw[transition, bend left=14, e-returnPaused-returningToChargingStation/.try] (s-returnPaused) to node[event, align=center] {Resume [ResumePermitted]} (s-returningToChargingStation);
\draw[transition, e-returnPaused-ready/.try] (s-returnPaused) to node[event, align=center] {Acknowledge} (s-ready);
\draw[transition, bend left=14, e-driving-avoidingObstacle/.try] (s-driving) to node[event, align=center] {ObstacleDetected} (s-avoidingObstacle);
\draw[transition, bend left=14, e-avoidingObstacle-driving/.try] (s-avoidingObstacle) to node[event, align=center] {ObstacleCleared} (s-driving);
\draw[transition, e-avoidingObstacle-blocked/.try] (s-avoidingObstacle) to node[event, align=center] {RouteBlocked} (s-blocked);
\draw[transition, e-paused-stopped/.try] (s-paused) to node[event, align=center] {Cancel} (s-stopped);
\draw[transition, e-stopped-ready/.try] (s-stopped) to node[event, align=center] {Acknowledge} (s-ready);
\draw[transition, e-completed-ready/.try] (s-completed) to node[event, align=center] {Acknowledge} (s-ready);
\draw[transition, bend left=14, e-rejected-ready/.try] (s-rejected) to node[event, align=center] {Acknowledge} (s-ready);
\draw[transition, e-blocked-ready/.try] (s-blocked) to node[event, align=center] {Acknowledge} (s-ready);
\draw[transition, e-chargingInterfaceReached-ready/.try] (s-chargingInterfaceReached) to node[event, align=center] {Acknowledge} (s-ready);
\draw[transition, e-fault-ready/.try] (s-fault) to node[event, align=center] {Acknowledge [faultCleared]} (s-ready);
\draw[transition, e-safeState-ready/.try] (s-safeState) to node[event, align=center] {Acknowledge [safeToMove]} (s-ready);
\node[note] at (pos-note) {\textbf{Sicherheit (DEC-016):} Aus jedem Zustand ausser \texttt{off}, \texttt{fault}, \texttt{safeState} fuehrt \texttt{UnsafeCondition} nach \texttt{safeState} und \texttt{TechnicalFault} nach \texttt{fault}. Diese 26 Uebergaenge sind nicht einzeln gezeichnet. Ereignisnamen ohne die Endung \texttt{Requested}.};
\end{tikzpicture}
